% =====================================================================
\chapter{Autoevaluación}
\label{ap:autoevaluacion}
% =====================================================================

\section{Cálculo de la nota}

La nota del examen es la suma de los puntos de la parte~A (máximo 30) y de la parte~B (máximo 70). El examen se aprueba con \textbf{70 puntos o más}. Si se obtienen menos de 15 puntos en la parte~A o menos de 40 en la parte~B, conviene repasar esa parte aunque el examen esté aprobado: la Fase~II (módulos 7 a~11) da por sabido todo este material.

\begin{center}
\begin{tabular}{@{}llcc@{}}
\toprule
Parte & Elemento & Puntos & Obtenidos \\
\midrule
A & Cuestiones 1--20 (1,5 cada una) & 30 & \\
\midrule
B & B.1 Modelo de bigramas & 12 & \\
B & B.2 Embeddings y analogías & 10 & \\
B & B.3 Autovalores, SVD y bajo rango & 12 & \\
B & B.4 Gradiente de la softmax con entropía cruzada & 14 & \\
B & B.5 Estimación y regla de Bayes & 11 & \\
B & B.6 Entropía, compresión y bits por byte & 11 & \\
\midrule
 & \textbf{Total} & \textbf{100} & \\
\bottomrule
\end{tabular}
\end{center}

\section{Puntuación por módulo}

La tabla siguiente agrupa los puntos por módulo. Un módulo en el que se obtenga menos del $70\,\%$ de sus puntos debe repasarse aunque el examen esté aprobado.

\begin{center}
\begin{tabular}{@{}lccccc@{}}
\toprule
Módulo & Cuestiones & Puntos A & Problema & Puntos B & Total \\
\midrule
M01 & 1--4 & 6 & B.1 & 12 & 18 \\
M02 & 5--6 & 3 & B.2 & 10 & 13 \\
M03 & 7--11 & 7,5 & B.3 & 12 & 19,5 \\
M04 & 12--14 & 4,5 & B.4 & 14 & 18,5 \\
M05 & 15--17 & 4,5 & B.5 & 11 & 15,5 \\
M06 & 18--20 & 4,5 & B.6 & 11 & 15,5 \\
\midrule
\textbf{Total} & 20 & \textbf{30} & 6 & \textbf{70} & \textbf{100} \\
\bottomrule
\end{tabular}
\end{center}

\section{Puerta de salida de la Fase I}

Para pasar a la Fase~II deben cumplirse las tres condiciones siguientes:
\begin{enumerate}[label=\arabic*.]
  \item un resultado de al menos 9 sobre 12 en el test de cada módulo, M01 a M06;
  \item una autoevaluación del laboratorio interactivo de la fase de al menos el $70\,\%$;
  \item una nota de este examen de hito de al menos el $70\,\%$ (70 puntos sobre 100).
\end{enumerate}
El proyecto de fase se califica con la rúbrica del capítulo correspondiente y cuenta en la nota de la fase.

\section{Qué repasar según los fallos}

\begin{center}
\small
\renewcommand{\arraystretch}{1.1}
\begin{tabularx}{\textwidth}{@{}>{\raggedright\arraybackslash}p{3.6cm}X@{}}
\toprule
Fallo & Secciones de los apuntes \\
\midrule
Cuestiones 1 y 2 & M01, §2.2 (regla de la cadena) y §2.3 ($n$-gramas) \\
Cuestiones 3 y 4 & M01, §5.1 (perplejidad) y §4.1 (softmax con temperatura) \\
B.1 (a), (b) y (e) & M01, §2.3, §3.2 y §4.3 \\
B.1 (c) y (d) & M01, §5.1, §5.3 y §3.3 (suavizado aditivo) \\
Cuestión 5 y B.2 (a) & M02, §1.2, §1.3 y §2.3 \\
Cuestión 6 y B.2 (b) & M02, §5.1 (matriz de embeddings) y §5.4 (analogías) \\
B.2 (c) y (d) & M02, §3.3 y §4.4; M03, §1.3 (núcleo) \\
Cuestiones 7 y 11 & M03, §6.2 (producto por lotes) y §6.3 (\emph{broadcasting}) \\
Cuestión 8 y B.3 (a) y (b) & M03, §1.3, §2.1, §3.1, §3.2 y §3.3 \\
Cuestión 9 y B.3 (c) y (d) & M03, §4.1, §4.2, §5.1, §5.2 y §5.3 \\
Cuestión 10 & M03, §5.4 (PCA) \\
Cuestión 12 & M04, §1.1 y §1.2 \\
Cuestión 13 y B.4 (a)--(c) & M04, §3.2, §3.3 y §3.4 \\
B.4 (d) & M04, §1.5 y §4.4 \\
Cuestión 14 y B.4 (e) & M04, §6.1 y §6.3 \\
Cuestión 15 & M05, §2.3 \\
Cuestión 16 y B.5 (a) & M05, §4.2; M01, §3.2 \\
B.5 (b) & M05, §1.3 \\
Cuestión 17 y B.5 (c) & M05, §5.4, §5.5 y §5.6 \\
Cuestión 18 y B.6 (a) & M06, §1.2, §3.2, §3.3 y §3.4 \\
Cuestiones 19 y 20 y B.6 (b) & M06, §2.2 y §2.3 \\
B.6 (c) & M06, §4.2, §4.3 y §4.4; M01, §5.1 \\
\bottomrule
\end{tabularx}
\end{center}

\section{Qué hacer si no se alcanza el umbral}

\begin{enumerate}[label=\arabic*.]
  \item Anotar cada fallo en el registro de errores, clasificado como error de \emph{concepto}, de \emph{cálculo} o \emph{despiste}.
  \item Para cada módulo por debajo del $70\,\%$, repetir el ciclo de estudio del método del curso con las secciones indicadas en la tabla anterior: teoría, laboratorio, práctica, ejercicios $\star\star\star$ del tema, test y tarjetas de repaso.
  \item Rehacer los problemas fallados sin mirar la solución y, a continuación, los ejercicios del mismo tema en los apuntes del módulo.
  \item Repetir el examen pasada al menos una semana, preferiblemente con datos cambiados, y no avanzar a la Fase~II hasta superar las tres condiciones de la puerta de salida.
\end{enumerate}
